% CEUR-WS abstract and keywords block.
% \maketitle must appear AFTER this block in samplepaper-ceur.tex.

\begin{abstract}
\textbf{Background.}~Large language models (LLMs) deliver strong
reasoning and generation performance, but their cost, latency, energy
demand, and hardware requirements remain high. Small language models
(SLMs) are better suited to edge and privacy-sensitive deployment,
yet their capabilities vary across domains.
\textbf{Method.}~This paper presents a systematic literature review
of 54 studies on collaborative SLM--LLM systems spanning hybrid,
multi-agent, retrieval-centered, and ensemble architectures, with
monolithic LLMs treated as baselines rather than as a primary
architecture class. Guided by four research questions on performance,
efficiency, orchestration, and domain suitability, we apply a PRISMA
2020-based review design with a three-pass qualitative coding scheme
that synthesizes models, architectures, datasets, metrics, latency,
cost, energy, and reported limitations.
\textbf{Results.}~Corpus analysis shows that 70\% of studies date to
2025 and hybrid architectures are the most prevalent design; outcome
coding reveals that strong results are reliably associated with
explicit control layers (routers, retrievers, planners) rather than
with model scale.
\textbf{Conclusions.}~SLMs are most competitive when paired with
routing, retrieval, planning, or fusion mechanisms in narrow,
structured, or domain-specific settings.
\end{abstract}

\begin{keywords}
systematic literature review \sep
large language models \sep
small language models \sep
hybrid architectures \sep
multi-agent systems \sep
edge-cloud inference
\end{keywords}
